% ============================================================
% Study 1 — Baseline overconfidence in 1-day equity predictions
% ============================================================

\subsection*{Study 1}

Study~1 established a baseline calibration estimate for three frontier LLMs
producing confidence intervals for 1-day equity price predictions. We predicted
that all models would be overconfident---stating intervals narrower than warranted
by realized outcomes---and that the degree of miscalibration would differ
systematically across models.

\subsubsection*{Methods}

Three commercially deployed LLMs were evaluated: Claude
(claude-sonnet-4-20250514; Anthropic), GPT-4 (gpt-4o; OpenAI), and Gemini
(gemini-2.5-flash; Google DeepMind). All models were queried via their public
APIs at default temperature with a maximum of 512 output tokens per call.

\paragraph{Stimuli and predictions} Twenty large-cap S\&P~500 equities
were selected as prediction targets (AAPL, NVDA, MSFT, AMZN, GOOGL, META,
AVGO, TSLA, BRK-B, WMT, LLY, JPM, XOM, V, JNJ, MU, MA, ORCL, COST, SPY).
Each prompt provided the current closing price and prediction horizon as plain
text, requesting a point estimate plus 50\%, 80\%, and 90\% confidence
intervals as structured JSON. No chain-of-thought reasoning or calibration
instructions were provided; models were evaluated in their default,
deployment-ready configuration. Predictions were collected on March~23, 2026
for outcomes on March~24, 2026. Each model made 5 independent runs per
ticker, yielding 100 scored predictions per model.

\paragraph{Calibration metrics} For each confidence level
$\alpha \in \{0.50, 0.80, 0.90\}$, the empirical \textit{hit rate}
$\hat{p}_\alpha$ is the proportion of predictions for which the true outcome
fell within the stated interval $[L, U]$. A calibrated model achieves
$\hat{p}_\alpha = \alpha$; values below indicate overconfidence. Expected
Calibration Error (ECE) is the mean absolute deviation between stated and
empirical coverage across all three levels:
\begin{equation}
  \label{eq:ece}
  \text{ECE} = \frac{1}{3} \sum_{\alpha \in \{0.50,\,0.80,\,0.90\}}
               \left| \alpha - \hat{p}_\alpha \right|.
\end{equation}
A perfectly calibrated model achieves ECE\,=\,0.

\subsubsection*{Results}

Table~\ref{tab:study1} summarizes calibration for all three models.

\paragraph{Hit rates} All models were substantially overconfident at every
confidence level. Claude was the least miscalibrated (Hit@80\,=\,58\%,
Hit@90\,=\,73\%), but remained far below calibration targets. GPT-4's 80\%
CI contained the true outcome only 31\% of the time, and its 90\% CI only
38\%---implying interval widths roughly consistent with a 30\% confidence
level. Gemini closely tracked Claude at the 80\% level but underperformed at
90\% (65\%).

\paragraph{ECE} GPT-4's overall ECE (43.7\%) was more than double Claude's
(19.7\%) and more than triple the perfect-calibration target of 0\%.
Gemini was intermediate (22.7\%).

\begin{table}[ht]
  \centering
  \caption{Study~1: Calibration in 1-day equity predictions ($n = 100$ per model).
    ECE = Expected Calibration Error; lower is better.}
  \label{tab:study1}
  \begin{tabular}{lcccc}
    \toprule
    Model   & Hit@50\% & Hit@80\% & Hit@90\% & ECE    \\
    \midrule
    Claude  & 30.0\%   & 58.0\%   & 73.0\%   & 19.7\% \\
    Gemini  & 29.0\%   & 58.0\%   & 65.0\%   & 22.7\% \\
    GPT-4   & 20.0\%   & 31.0\%   & 38.0\%   & 43.7\% \\
    \midrule
    Perfect & 50.0\%   & 80.0\%   & 90.0\%   & 0.0\%  \\
    \bottomrule
  \end{tabular}
\end{table}

\subsubsection*{Discussion}

All three models produced systematically narrow confidence intervals for a
simple, single-domain forecasting task. GPT-4's calibration (31\% coverage at
the 80\% level) is comparable to the most severe human expert overconfidence
on record: \citet{BenDavid2013} found that CFOs' 80\% intervals for annual
equity returns achieved only 36--38\% coverage. These results confirm
universal overconfidence in the simplest possible setting and establish the
model ranking---Claude best, GPT-4 worst, Gemini intermediate---that will
prove consistent across Studies~2 and~3 in the aggregate.
